\chapter{Follicle-Stimulating Hormone (FSH)}

\section*{What Is This Test?}
Follicle-stimulating hormone (FSH) is a glycoprotein hormone produced by gonadotrope cells of the anterior pituitary gland. It is composed of a common alpha subunit (shared with LH, TSH, and hCG) and a unique beta subunit that confers receptor specificity. FSH is secreted in a pulsatile manner under the influence of gonadotropin-releasing hormone (GnRH) from the hypothalamus. In females, FSH stimulates ovarian follicular growth and maturation, promotes granulosa cell proliferation, and induces aromatase activity to convert androgens to estradiol. FSH also stimulates inhibin B production by granulosa cells, which feeds back to selectively suppress FSH secretion. In males, FSH acts on Sertoli cells to stimulate spermatogenesis and inhibin B production. FSH levels vary across the menstrual cycle with a mid-cycle surge alongside LH. The day 3 FSH level (early follicular phase) is the standard marker for ovarian reserve assessment. FSH measurement is essential for evaluating ovarian function, menopause, pubertal development, male and female infertility, and classifying hypogonadism as primary (gonadal) or secondary (pituitary/hypothalamic).

\section*{Alternative Names}
Follitropin; Serum FSH; Plasma FSH; Day 3 FSH; FSH/LH Ratio; Gonadotropin, Follicle-Stimulating; Follicular Stimulating Hormone

\section*{Principle}
FSH is measured using a two-site sandwich immunometric (chemiluminescent) immunoassay on automated analyzers. Two monoclonal antibodies specific for the beta subunit of FSH (with minimal cross-reactivity to LH, TSH, or hCG) are used: a capture antibody bound to a solid phase (magnetic particles or microtiter wells) and a detection antibody labeled with a chemiluminescent tag (ruthenium on Roche ECLIA, acridinium ester on Abbott Architect/Siemens). After incubation and washing, the chemiluminescent signal is proportional to FSH concentration. Current third-generation chemiluminescent assays have functional sensitivity of <0.1 mIU/mL and are calibrated against the WHO 2nd International Standard (IS 94/632) \cite{tietz2023textbook}. Inter-assay precision (coefficient of variation) is typically less than 5\% at clinically relevant concentrations. Because FSH is a glycoprotein with heterogeneous glycosylation patterns (differing sialic acid content), assays from different manufacturers may produce slightly different absolute values. The same assay and laboratory should be used for serial monitoring. LC-MS/MS is not used for FSH due to its glycoprotein structure and low circulating concentrations.

\section*{History}
FSH was identified as a distinct gonadotropin from LH in the 1930s through pituitary extraction and fractionation studies. In 1957, the first FSH bioassay (the Steelman-Pohley assay, measuring ovarian weight augmentation in hCG-treated rats) was developed. The first radioimmunoassay for FSH was developed by Rees Midgley in 1967 at the University of Michigan \cite{midgley1967radioimmunoassay}. The purification of human menopausal gonadotropin (hMG) from postmenopausal urine in the 1960s made ovulation induction clinically possible. Recombinant FSH (follitropin alfa, Gonal-F, marketed 1995; follitropin beta, Follistim, marketed 1997) was developed using Chinese hamster ovary (CHO) cell expression systems, revolutionizing infertility treatment by providing a pure, unlimited supply of FSH. The concept of using day 3 FSH as an ovarian reserve test was pioneered by Richard Scott and colleagues in the late 1980s at the Jones Institute. Anti-Mullerian hormone (AMH) has emerged as a superior marker of ovarian reserve since the 2010s, though FSH remains widely used.

\section*{Why This Test Is Ordered}
\begin{itemize}
  \item Assessment of ovarian reserve in infertility evaluation --- day 2--4 (early follicular phase) FSH >10--15 mIU/mL indicates diminished ovarian reserve and predicts poor response to ovarian stimulation during IVF \cite{asrm2020testing}
  \item Diagnosis of menopause (FSH >30--40 mIU/mL with amenorrhea >12 months in women >45 years) or premature ovarian insufficiency (FSH >40 mIU/mL on 2 measurements >1 month apart in women <40 years) \cite{webber2016eshre}
  \item Evaluation of primary amenorrhea (elevated FSH suggests gonadal dysgenesis/Turner syndrome, 45,X; low FSH suggests hypogonadotropic hypogonadism/Kallmann syndrome)
  \item Evaluation of secondary amenorrhea and oligomenorrhea --- distinguish ovarian failure from hypothalamic/pituitary causes
  \item Diagnosis of polycystic ovary syndrome (PCOS) --- LH:FSH ratio >2:1 is supportive but not diagnostic per Rotterdam criteria \cite{rotterdam2004revised}
  \item Classification of male hypogonadism: elevated FSH with low testosterone indicates primary testicular failure (Klinefelter syndrome 47,XXY, Sertoli cell dysfunction); low or inappropriately normal FSH with low testosterone indicates secondary hypogonadism
  \item Evaluation of male infertility --- elevated FSH with azoospermia/oligospermia suggests impaired spermatogenesis (non-obstructive); normal FSH with azoospermia suggests obstruction
  \item Assessment of central precocious puberty --- elevated FSH and LH in children <8 years (girls) or <9 years (boys) with secondary sexual characteristics
  \item Monitoring response to GnRH agonist/antagonist therapy in assisted reproduction and prostate cancer
  \item Investigation of pituitary function in suspected hypopituitarism --- low FSH/LH with low estradiol or testosterone suggests gonadotropin deficiency
\end{itemize}

\section*{Is This a Primary or Secondary Test}
FSH is a primary test for evaluating ovarian failure, menopause, and classifying male/female hypogonadism. It is part of the standard infertility evaluation in women (day 3 FSH with estradiol and AMH) and hypogonadism evaluation in men (FSH, LH, total and free testosterone). In PCOS, it is a secondary (supportive) test.

\section*{How This Test Is Performed}
A healthcare professional draws a blood sample from a vein into a serum separator tube (gold-top) or plain red-top tube, collecting approximately 3--5 mL of blood. For ovarian reserve testing in menstruating women, the sample MUST be drawn on day 2, 3, or 4 of the menstrual cycle (early follicular phase), when estradiol and inhibin B are at their lowest and FSH is maximally released from negative feedback. The cycle day must be documented on the laboratory requisition. For menopause evaluation in amenorrheic women, FSH can be drawn on any day. In men, FSH can be drawn at any time, though morning (7--10 AM) is preferred for concurrent testosterone measurement. FSH is relatively stable at room temperature for 24 hours and refrigerated for 7 days.

\section*{What Preparation Is Needed}
No fasting is required. In menstruating women, the cycle day must be documented. Estrogen therapy (oral contraceptives, hormone replacement therapy) suppresses FSH; ideally, these should be held for 4--6 weeks before testing if evaluating pituitary-gonadal function. GnRH agonists and antagonists (used in IVF and prostate cancer) profoundly suppress or stimulate FSH; results must be interpreted in this context. High-dose biotin (>5 mg/day) should be discontinued 48 hours before testing due to streptavidin-biotin assay interference. Anabolic steroids, exogenous testosterone, and danazol suppress FSH. Clomiphene citrate and tamoxifen (SERMs) elevate FSH by blocking estrogen negative feedback on the hypothalamus.

\section*{Contraindications and Precautions}
FSH values MUST be interpreted in the context of menstrual cycle phase, patient age, sex, and concurrent medications. A single FSH in a menstruating woman drawn at an unknown cycle phase is largely uninterpretable. Oral contraceptives and GnRH agonists/antagonists profoundly suppress FSH. In the perimenopausal transition, FSH fluctuates widely from month to month and is not reliable. Day 3 FSH is only interpretable when day 3 estradiol is also low (<80 pg/mL); an elevated estradiol with a normal FSH may be falsely reassuring because estradiol suppresses FSH. Chronic illness, malnutrition, and severe stress (functional hypothalamic amenorrhea) suppress FSH. In men, FSH increases mildly with normal aging (sertoli cell function declines). In chronic kidney disease, FSH may be mildly elevated due to decreased renal clearance. Acute illness transiently suppresses FSH through central mechanisms. Pulsatile secretion causes 20--30\% variation in FSH within hours --- a single value may not reflect the integrated secretion, though FSH is less pulsatile than LH.

\section*{Risks}
Minimal risk. Slight pain, bruising, or hematoma at venipuncture site. Rarely: vasovagal syncope or infection.

\section*{What Happens After the Test}
FSH results are available within 1--4 hours on automated analyzers. In women: elevated day 3 FSH (>10--15 mIU/mL) indicates diminished ovarian reserve. FSH >30--40 mIU/mL on 2 measurements >1 month apart with amenorrhea in a woman <40 years indicates premature ovarian insufficiency --- further workup includes karyotype (Turner syndrome), FMR1 premutation testing (fragile X-associated POI), and anti-adrenal antibodies (autoimmune polyglandular syndrome). FSH >40 mIU/mL with amenorrhea >12 months in a woman >45 years confirms menopause. In men: elevated FSH with low testosterone indicates primary testicular failure --- karyotype for Klinefelter syndrome (47,XXY) is indicated. High FSH with azoospermia suggests non-obstructive etiology (spermatogenic failure). Normal/low FSH with azoospermia suggests obstruction (congenital bilateral absence of the vas deferens). Low FSH with low testosterone indicates secondary hypogonadism, requiring pituitary MRI and evaluation of other pituitary hormones.

\section*{Understanding Results}

\subsection*{Below Normal}
\begin{itemize}
  \item Secondary (hypogonadotropic) hypogonadism: low or inappropriately normal FSH with low estradiol (female) or low testosterone (male). Causes include functional hypothalamic amenorrhea (excessive exercise, low body weight, eating disorders, stress), Kallmann syndrome (congenital GnRH deficiency with anosmia/hyposmia), idiopathic hypogonadotropic hypogonadism, pituitary adenomas (prolactinoma, non-functioning macroadenoma compressing gonadotropes), pituitary surgery or radiation, Sheehan syndrome, hemochromatosis, and chronic opioid use suppressing GnRH.
  \item PCOS: FSH is typically normal or low-normal (unlike LH which is elevated), resulting in an elevated LH:FSH ratio >2:1 seen in approximately 60\% of PCOS patients.
  \item Exogenous estrogen or testosterone therapy: negative feedback suppresses FSH. Oral contraceptives typically suppress FSH to <5 mIU/mL.
  \item Hyperprolactinemia: elevated prolactin suppresses GnRH, reducing FSH and LH secretion.
  \item Constitutional delay of growth and puberty: temporary delay in pubertal onset with low gonadotropins for chronologic age but appropriate for bone age.
\end{itemize}

\subsection*{Above Normal}
\begin{itemize}
  \item Menopause: FSH >30--40 mIU/mL (often >100 mIU/mL) with amenorrhea >12 months in women >45 years. Ovarian follicle depletion eliminates inhibin B and estradiol negative feedback, causing unopposed FSH secretion. FSH rises before LH and reaches higher absolute levels.
  \item Premature ovarian insufficiency (POI): FSH >40 mIU/mL on 2 measurements >1 month apart in women <40 years. Affects 1\% of women. Causes: genetic (Turner syndrome 45,X, FMR1 premutation in 3--5\%); autoimmune (APS-1, APS-2, isolated autoimmune oophoritis); iatrogenic (chemotherapy, especially alkylating agents like cyclophosphamide; pelvic radiation); metabolic (galactosemia); idiopathic (majority).
  \item Primary testicular failure: elevated FSH (often >15--20 mIU/mL, may exceed 50 mIU/mL) with low or low-normal testosterone. Klinefelter syndrome (47,XXY) is the most common cause --- small firm testes (<4 mL), gynecomastia, azoospermia, tall stature with eunuchoid proportions. Other causes: cryptorchidism, orchitis (mumps), testicular torsion, chemotherapy/radiation, and Sertoli-cell-only syndrome.
  \item Diminished ovarian reserve: day 3 FSH >10--15 mIU/mL in women <40 years with regular cycles. Predicts reduced response to ovarian stimulation and lower IVF success rates. AMH and antral follicle count (AFC) are more sensitive markers.
  \item Gonadal dysgenesis: elevated FSH in infants/children with Turner syndrome or Swyer syndrome (46,XY complete gonadal dysgenesis).
  \item Azoospermia with elevated FSH: non-obstructive azoospermia (spermatogenic failure). FSH is typically >2--3 times upper limit of normal. Testicular biopsy may show maturation arrest, hypospermatogenesis, or Sertoli-cell-only pattern.
  \item Physiologic FSH elevation: perimenopausal transition (marked FSH variability, levels fluctuate 20--150 mIU/mL from month to month), and the mid-cycle FSH surge (FSH peaks at 10--20 mIU/mL alongside LH surge >25--50 mIU/mL).
\end{itemize}

\section*{Normal Results}
\begin{itemize}
  \item Adult female, follicular phase (day 2--4): 2.5--10 mIU/mL
  \item Adult female, mid-cycle surge: 5--20 mIU/mL
  \item Adult female, luteal phase: 1.5--5 mIU/mL
  \item Adult female, postmenopausal: 25--135 mIU/mL (may reach >200 mIU/mL)
  \item Adult male (18--60 years): 1.5--12 mIU/mL
  \item Adult male (>60 years): 2--15 mIU/mL (mild age-related increase)
  \item Children (pre-pubertal): <3 mIU/mL
  \item Ovarian reserve cutoff (day 3 FSH): <10 mIU/mL (normal); 10--15 (diminished, intermediate prognosis); >15 (significantly diminished, poor prognosis for IVF)
  \item Premature ovarian insufficiency cutoff: FSH >40 mIU/mL on 2 measurements >1 month apart
  \item Note: Reference ranges vary by laboratory and assay. Always use laboratory-specific ranges. FSH values should be interpreted with estradiol and LH.
\end{itemize}

\section*{Follow-up Tests}
\begin{itemize}
  \item LH (luteinizing hormone): Always co-measured with FSH to assess the LH:FSH ratio and classify gonadal axis disorders. LH:FSH >2:1 supports PCOS.
  \item Estradiol (day 3): Must be measured together with day 3 FSH. Estradiol should be <80 pg/mL on day 3 for FSH to be interpretable.
  \item Total and free testosterone (male): Classify primary vs. secondary hypogonadism. Morning sample (7--10 AM) \cite{bhasin2018testosterone}.
  \item AMH (anti-Mullerian hormone): Superior marker of ovarian reserve, correlates with antral follicle count, does not require cycle day timing.
  \item Inhibin B: Direct marker of granulosa cell (female) and Sertoli cell (male) function \cite{burger2018physiology}. Low inhibin B supports diminished ovarian reserve or Sertoli cell dysfunction.
  \item Prolactin: Exclude hyperprolactinemia as a cause of hypogonadotropic hypogonadism.
  \item TSH and free T4: Exclude thyroid dysfunction causing menstrual disturbances.
  \item Karyotype (chromosome analysis): Indicated for primary testicular failure (Klinefelter syndrome, 47,XXY) and premature ovarian insufficiency (Turner syndrome, 45,X; X chromosome deletions).
  \item FMR1 (fragile X) premutation testing: In women with premature ovarian insufficiency (3--5\% carry FMR1 premutation).
  \item Pituitary MRI with gadolinium: For secondary hypogonadism (low FSH/LH with low estradiol or testosterone) or visual field defects.
  \item Semen analysis: In male infertility evaluations. Azoospermia with normal FSH suggests obstruction; azoospermia with elevated FSH suggests non-obstructive etiology.
\end{itemize}

\section*{References}
\bibliographystyle{unsrtnat}
\bibliography{references}

\clearpage
